% id: 125-28
% book: 베이직쎈 확률과 통계 · 07 통계적 추정
% section: 개념 38 표본평균의 평균, 분산, 표준편차
% passage: 모집단의 확률변수 $X$의 확률분포를 표로 나타내면 아래와 같다. 이 모집단에서 임의추출한 크기가 9인 표본의 표본평균을 $\overline{X}$라 할 때, 다음을 구하시오. \cond{$a$는 상수이다.}
% figure: none
% answer: 
% answer_source: 
% variant_level: 0 (원본 전사)
% 이 파일은 build.mjs 가 items.json 에서 만든다. 고칠 때는 items.json 을 고친다.
\smash{\raisebox{1.5mm}{\dmkichul{베이직쎈 확률과 통계 125쪽 28번}}}
\par\smallskip\dispeq{{\setlength{\tabcolsep}{7pt}\renewcommand{\arraystretch}{2.2}\begin{tabular}{|c|c|c|c|c|}\hline $X$ & $1$ & $3$ & $5$ & 합계 \\ \hline $\mathrm{P}(X=x)$ & $a$ & $\dfrac{1}{5}$ & $\dfrac{2}{5}$ & $1$ \\ \hline\end{tabular}}}
\dmsubprob{1} $a$의 값
\dmsubprob{2} $\mathrm{E}(\overline{X})$, $\mathrm{V}(\overline{X})$
